\section{Discussion}
\label{sec:limits}
R20 produced six complete task-to-approval-to-receipt-to-sample traces across
two simulator scenes; none of its single post-readback samples reached target.
R21 passed 31 assertions on the same configured route, and a 12-case mock
suite refused each of four stale approval-context mutations before actuation.
Together, these results verify selected software behavior, not manufacturing
performance.

The runs do not establish independent replication, process convergence, or
quality improvement. Each scene used one simulator instance, and its single
probe is not a response trajectory. Configuration reuse, alternate harnesses
or direct REST entry, online DAQ, operator effectiveness, physical transfer,
and comparative architecture value remain untested. Synthetic approvals do
not assess operator effectiveness
\cite{parasuramanmanzey2010complacency,parasuraman2000model,bainbridge1983ironies};
restricted tool access does not establish deployment security
\cite{greshake2023injection,liu2024pinjection,iec62443-3-3} or replace
independently engineered PLC/SIS protection \cite{iec61511}. The
context-mutation suite does not compare architectures or quantify integration
effort. Future work should use a defined baseline, exercise alternate routes
and online DAQ, and include independent simulator and physical-plant trials.

\section{Conclusion}
\label{sec:conclusion}
AgentWorkShop connects collaborative tasks to production context, checked
operations, semantic I/O, and service-owned records. R20 and R21 verify
selected paths across two simulator mappings; a mock-state suite confirms
refusal when an approved line--product--recipe--run tuple changes. The evidence
supports the implemented software contract, not process-control efficacy,
comparative advantage, or portability.
